%! TEX root = ../am2-20-21.tex
\chapter{Integrale de suprafață}

\section{Integrale de suprafață de speța întîi}
\index{integrale!de suprafață!speța întîi}
Integralele de suprafață reprezintă generalizarea într-o dimensiune
superioară pentru integralele curbilinii. Astfel, multe dintre formulele
și abordările de calcul pe care le vom folosi vor fi similare.

Fie $ \Phi : D \to \RR^3 $ o pînză parametrizată și fie $ \Sigma = \Phi(D) $
imaginea ei (i.e.\ suprafață pe care vom integra). Fie $ f : U \seq \Sigma \to \RR $
o funcție continuă (deci integrabilă), definită pe (o porțiune din) imaginea
pînzei.

Vom fi interesați de un caz particular (și, totodată, cel mai des întîlnit) pentru
integrala de suprafață, anume cînd pînza este dată într-o \emph{parametrizare %
  carteziană}. Adică ecuația suprafeței poate fi scrisă în forma $ z = z(x, y) $.

Integrala de suprafață de speța întîi a funcției $ f $ pe suprafața $ \Sigma $
parametrizată cartezian prin $ z = z(x, y) $ este:
\[
  \int_\Sigma f(x, y, z) d\sigma = \iint_D f(x, y, z(x, y)) \cdot %
  \sqrt{1 + p^2 + q^2} dxdy,
\]
unde:
\begin{itemize}
\item $ D $ este domeniul de definiție al parametrizării carteziene,
  adică proiecția pe planul XOY a suprafeței ($ z : D \seq \RR^2 \to \RR $);
\item coeficienții de sub radical sînt $ p = \dfrac{\partial z}{\partial x} $
  și $ q = \dfrac{\partial z}{\partial y} $.
\end{itemize}

În cazul particular în care $ f = 1 $, se obține \emph{aria suprafeței $ \Sigma $}.

%%%%%%%%%%%%%%%%%%%%%%%%%%%%%%%%%%%%%%%%%%%%%%%%%%%%%%%%%%%%%%%%%%%%%% 

\section{Integrale de suprafață de speța a doua}
\index{integrale!de suprafață!speța a doua}
Fie, ca mai sus, o pînză tridimensională, pe care o considerăm a fi parametrizată:
\[
  \Phi : D \to \RR^3, \quad \Phi(u, v) = (X(u, v), Y(u, v), Z(u, v)).
\]

Fie, de asemenea, o 2-formă diferențială\footnotemark:
\[
  \omega = P dy\wedge dz + Q dz \wedge dx + R dx \wedge dy.
\]

\footnotetext{$ \omega $ este o \textbf{2}-formă diferențială deoarece
  ea conține produse de cîte 2 elemente diferențiale $ dx, dy, dz $. De asemenea,
  notația $ \wedge $ (citită \qq{wedge} sau \qq{produs exterior}) este o
  notație specifică pentru produsul care se definește între diferențialele
  $ dx, dy, dz $. Alternativ, puteți găsi scrierea și prin juxtapunere:
  \[
    \omega = P dydz + Q dzdx + R dxdy.
  \]
  Însă acest produs \emph{nu este comutativ}, deci ordinea în care scriem
  2-forma diferențială este esențială! (observați permutările circulare:
  $ dydz \to dzdx \to dxdy $).}

Integrala pe suprafața orientată $ \Sigma = \Phi(D) $ a 2-formei diferențiale
$ \omega $ se definește prin:
\[
  \int_\Sigma \omega = \iint_D (P \circ \Phi) \cdot \frac{D(Y,Z)}{D(u, v)} + %
  (Q \circ \Phi) \cdot \frac{D(Z, X)}{D(u, v)} + %
  (R \circ \Phi) \cdot \frac{D(X, Y)}{D(u, v)} dudv,
\]
unde $ D $ este domeniul parametrizării ($(u, v) \in D$),
iar $ \dfrac{D(Y, Z)}{D(u, v)} $ etc.\ sînt jacobienii parametrizării
$ (X, Y, Z) $ în funcție de $ u, v $. Concret, de exemplu, avem:
\[
  \frac{D(X, Y)}{D(u, v)} = %
  \begin{vmatrix}
    X_u & X_v \\
    Y_u & Y_v
  \end{vmatrix},
\]
și celelalte, folosind notația simplificată $ X_u = \dfrac{\partial X}{\partial u} $ etc.
\footnote{Ordinea scrierii jacobienilor este, evident, esențială. Pentru
  a-i reține mai ușor, observați că ordinea urmează tot o permutare circulară,
  asemănătoare diferențialelor din 2-forma diferențială $ \omega $.}

Într-o formă simplificată, putem scrie integrala folosind un determinant formal:
\[
  \int_\Sigma \omega = \int %
  \begin{vmatrix}
    P \circ \Phi & Q \circ \Phi & R \circ \Phi \\
    X_u & Y_u & Z_u \\
    X_v & Y_v & Z_v
  \end{vmatrix} dudv.
\]

%%%%%%%%%%%%%%%%%%%%%%%%%%%%%%%%%%%%%%%%%%%%%%%%%%%%%%%%%%%%%%%%%%%%%%

\section{Formula Gauss-Ostrogradski}
\index{formula!Gauss-Ostrogradski}
\index{formula!flux-divergență}
Această formulă ne permite să schimbăm o integrală de suprafață de speța
a doua cu una triplă, similar formulei Green-Riemann, dar în dimensiune
superioară.

Păstrînd notațiile și contextul de mai sus, formula Gauss-Ostrogradski se
scrie:
\[
  \int_\Sigma \omega = \iiint_K P_x + Q_y + R_z dxdydz,
\]
unde $ K = \dr{Int}\Sigma $ este solidul care are drept frontieră suprafața
$ \Sigma $, iar $ P_x, Q_y, R_z $ notează derivatele parțiale corespunzătoare
coeficienților din 2-forma diferențială $ \omega $.

\begin{remark}\label{rk:go-exist}
  Formula Gauss-Ostrogradski are, ca și formula Green-Riemann, \emph{condiții %
  de aplicare (existență)}. Încercați să le formulați, analizînd formula.
\end{remark}

Formula Gauss-Ostrogradski se mai numește \emph{formula flux-divergență}.
Într-adevăr, folosind o interpretare fizică, se poate asocia 2-formei diferențiale
$ \omega $ cîmpul vectorial $ \vec{V} = (P, Q, R) $, iar membrul stîng, adică
integrala de suprafață, calculează \emph{fluxul cîmpului $ \vec{V}$ prin suprafața
  $ \Sigma $}. În fizică, acesta se definește ca produsul scalar dintre cîmpul
vectorial și versorul normal la suprafață. Membrul drept este, după cum se poate
vedea ușor, divergența cîmpului vectorial \emph{în solidul $ \dr{Int}\Sigma $},
deci avem:
\[
  \int_\Sigma \vec{V} \cdot \vec{n} d\sigma = \iiint_K \nabla \cdot \vec{V} dxdydz.
\]

Pentru cazul cînd suprafața este parametrizată, adică avem:
\[
  \Phi = \Phi(X(u, v), Y(u, v), Z(u, v)),
\]
se pot calcula \emph{vectorii tangenți} la suprafață, după direcțiile
lui $ u $ și $ v $, prin derivate parțiale:
\[
  \vec{T}_u = \dfrac{\partial \Phi}{\partial u}, \qquad %
  \vec{T}_v = \dfrac{\partial \Phi}{\partial v}.
\]

Apoi, normala la suprafață se poate calcula prin produs vectorial\footnote{
  amintiți-vă, produsul vectorial al doi vectori este un vector perpendicular
  pe planul dat de cei doi factori}. O variantă simplă de a reține formula
de calcul pentru produsul vectorial folosește determinantul formal:
\[
  \vec{N} = \vec{T}_u \times \vec{T}_v = %
  \begin{vmatrix}
    \vec{i} & \vec{j} & \vec{k} \\
    \vec{T}_u^1 & \vec{T}_u^2 & \vec{T}_u^3 \\
    \vec{T}_v^1 & \vec{T}_v^2 & \vec{T}_v^3
  \end{vmatrix},
\]
unde $ \vec{T}_{u,v}^i $ notează componenta $ i $ a vectorului $ \vec{T}_{u, v} $.

Ulterior, \emph{vectorul} normal $ \vec{N} $ devine \emph{versorul} normal $ \vec{n} $,
prin normare, adică $ \vec{n} = \dfrac{1}{||\vec{N}||} \vec{N} $.

Dar, ținînd cont că avem relația între diferențiale $ d\sigma = ||\vec{N}|| dudv $,
rezultă că obținem:
\[
  \int_\Sigma \vec{V} \cdot \vec{n} d\sigma = \int_\Sigma \vec{V} \cdot \dfrac{1}{||\vec{N}||} \vec{N} \cdot ||\vec{N}|| dudv = %
  \iint_D \vec{V} \cdot \vec{N} dudv.
\]


%%%%%%%%%%%%%%%%%%%%%%%%%%%%%%%%%%%%%%%%%%%%%%%%%%%%%%%%%%%%%%%%%%%%%%

\section{Parametrizări uzuale}

Următoarele formule de parametrizare pot fi folosite în calcule:
\begin{enumerate}[(1)]
\item \textbf{Sfera:} Fie $ R > 0 $ și $ (u, v) \in D = [0, \pi] \times [0, 2\pi) $.
  Parametrizarea sferei $ \Phi = \Phi(u, v) $ este:
  \[
    \Phi(u, v) = (R \sin u \cos v, R \sin u \sin v, R \cos u);
  \]
\item \textbf{Elipsoidul:} Fie $ a, b, c > 0 $ și
  $ (u, v) \in D = [0, \pi] \times [0, 2\pi) $. Parametrizarea elipsoidului este:
  \[
    \Phi(u, v) = (a \sin u \cos v, b \sin u \sin v, c \cos u);
  \]
\item \textbf{Paraboloidul:} Fie $ a > 0, h > 0 $ și
  $ (u, v) \in D = [0, h] \times [0, 2\pi) $. Parametrizarea paraboloidului este:
  \[
    \Phi(u, v) = (au\cos v, au \sin v, u^2);
  \]
\item \textbf{Conul:} Fie $ h > 0 $ și $ (u, v) \in D = [0, 2\pi) \times [0, h] $.
  Parametrizarea conului este:
  \[
    \Phi(u, v) = (v \cos u, v \sin u, v);
  \]
\item \textbf{Cilindrul:} Fie $ a > 0, 0 \leq h_1 \leq h_2 $ și
  $ (u, v) \in [0, 2\pi] \times [h_1, h_2] $. Parametrizarea cilindrului
  este:
  \[
    \Phi(u, v) = (a \cos u, a \sin u, v).
  \]
\end{enumerate}


%%%%%%%%%%%%%%%%%%%%%%%%%%%%%%%%%%%%%%%%%%%%%%%%%%%%%%%%%%%%%%%%%%%%%%

\section{Exerciții}

\indent\indent 1. Calculați vectorii tangenți și versorul normalei la suprafețele
parametrizate din secțiunea anterioară.

\vspace{0.5cm}

2. Să se calculeze integrala de suprafață de prima speță:
\[
  \int_\Sigma f(x, y, z) d\sigma,
\]
unde $ f(x, y, z) = y \sqrt{z} $, iar $ \Sigma: x^2 + y^2 = 6z, z \in [0, 2] $.

\vspace{0.5cm}

3. Folosind integrala de suprafață, calculați aria suprafeței $ \Sigma $, unde:
\[
  \Sigma : 2z = 4 - x^2 - y^2, \quad z \in [0, 1].
\]

\vspace{0.5cm}

4. Calculați fluxul cîmpului vectorial $ \vec{V} $ prin suprafața $ \Sigma $
pentru:
\[
  \vec{V} = y\vec{i} + x\vec{j} + z^2 \vec{k}, \quad %
  \Sigma : z = x^2 + y^2, z \in [0, 1].
\]

\vspace{0.5cm}

5. Calculați $ \ds\int_\Sigma \omega $, unde:
\begin{enumerate}[(a)]
\item
  \begin{itemize}
  \item $ \omega = y dy \wedge dz + z dz \wedge dx + x dx \wedge dy $;
  \item $ \Sigma : x^2 + y^2 = z^2, z \in [1, 2] $.
  \end{itemize}
\item
  \begin{itemize}
  \item $ \omega = x(z + 3) dy \wedge dz + yz dz \wedge dx - %
    (z + z^2) dx \wedge dy $;
  \item $ \Sigma : x^2 + y^2 + z^2 = 1, z \geq 0 $.
  \end{itemize}
\end{enumerate}

\vspace{0.5cm}

6. Calculați fluxul cîmpului:
\[
  \vec{V} = (x + y) dy \wedge dz + (y + z) dz \wedge dx -2z dx \wedge dy
\]
prin emisfera $ \Sigma: x^2 + y^2 + z^2 = 4, z > 0 $.

%%%%%%%%%%%%%%%%%%%%%%%%%%%%%%%%%%%%%%%%%%%%%%%%%%%%%%%%%%%%%%%%%%%%%%

\section{Formula lui Stokes}
\index{formula!Stokes}
Această ultimă formulă integrală ne permite să schimbăm o integrală curbilinie
de speța a doua cu una de suprafață de speța a doua.

Fie $ \Sigma $ o suprafață cu bord (frontieră) și fie o 1-formă diferențială
\[
  \alpha = Pdx + Qdy + Rdz,
\]
care este de clasă $ \kal{C}^1 $ într-o vecinătate a lui $ \Sigma $.

Notăm frontiera lui $ \Sigma $ prin $ \partial \Sigma $, care este o
curbă (conturul suprafeței) sau, mai precis, un \emph{drum neted}.

Are loc \emph{formula lui Stokes:}
\[
  \int_{\partial \Sigma} Pdx + Qdy + Rdz = %
  \int_\Sigma \Big( \frac{\partial R}{\partial y} - %
  \frac{\partial Q}{\partial z} \Big) dy \wedge dz + %
  \Big( \frac{\partial P}{\partial z} - %
  \frac{\partial R}{\partial x} \Big) dz \wedge dx + %
  \Big( \frac{\partial Q}{\partial x} - %
  \frac{\partial P}{\partial y}\Big) dx \wedge dy.
\]

În notație vectorială, dacă asociem cîmpul vectorial $ \vec{V} = (P, Q, R) $
1-formei diferențiale $ \alpha $, atunci formula lui Stokes se scrie:
\[
  \int_{\partial \Sigma} \vec{V} \cdot d\vec{r} = %
  \int_\Sigma (\nabla \times \vec{V}) \cdot \vec{n} d \sigma,
\]
unde $ \nabla \times \vec{V} = \dr{rot} \vec{V} $ se numește \emph{rotorul}
cîmpului vectorial, calculat cu ajutorul produsului vectorial formal între
operatorul diferențial $ \nabla = \left( \dfrac{\partial}{\partial x}, %
\dfrac{\partial}{\partial y}, \dfrac{\partial}{\partial z} \right) $ și
cîmpul $ \vec{V} = (P, Q, R) $.


%%%%%%%%%%%%%%%%%%%%%%%%%%%%%%%%%%%%%%%%%%%%%%%%%%%%%%%%%%%%%%%%%%%%%%

\section{Exerciții}

7. Să se calculeze, folosind formula lui Stokes, integrala curbilinie
$ \ds\int_\gamma \alpha $ pentru cazurile:
\begin{enumerate}[(a)]
\item $ \alpha = (y - z) dx + (z - x) dy + (x - y) dz $, iar curba $ \gamma $
  este frontiera suprafeței $ \Sigma : z = x^2 + y^2, z = 1 $;
\item $ \alpha = y dx + zdy + xdz $, iar curba $ \gamma $ este frontiera
  suprafeței $ \Sigma: x^2 + y^2 + z^2 = 1, x + y + z = 0 $.
\end{enumerate}

%%%%%%%%%%%%%%%%%%%%%%%%%%%%%%%%%%%%%%%%%%%%%%%%%%%%%%%%%%%%%%%%%%%%%%

\section{Resurse suplimentare}

\textbf{Formula Green-Riemann, formula Gauss-Ostrogradski} și \textbf{formula Stokes}
se numesc, în general, \emph{formule integrale}. Pentru coerență,
le-am introdus în capitolele potrivite, în loc să le dedic un capitol separat.

Recomand materialele Prof.\ Purtan, care conțin exerciții rezolvate complet:
\begin{itemize}
    \item \href{https://adrianmanea.xyz/docs/18-19-etti-am1/purtan-int-sup.pdf}{integrale de suprafață};
    \item \href{https://adrianmanea.xyz/docs/18-19-etti-am1/purtan-f-int.pdf}{formule integrale}.
\end{itemize}
